% Generated draft: requires mathematical and source review.
\documentclass[11pt]{article}
\usepackage[margin=25mm]{geometry}
\usepackage{fontspec}
\setmainfont{Latin Modern Roman}
\usepackage{amsmath,amssymb,mathtools,unicode-math}
\setmathfont{Latin Modern Math}
\usepackage{enumitem,needspace,xurl,hyperref}
\setlength{\parindent}{0pt}
\setlength{\parskip}{5pt}
\setlength{\emergencystretch}{3em}
\newcommand{\N}{\mathbb{N}}
\newcommand{\Z}{\mathbb{Z}}
\newcommand{\Q}{\mathbb{Q}}
\newcommand{\R}{\mathbb{R}}
\newcommand{\C}{\mathbb{C}}
\newcommand{\F}{\mathbb{F}}
\newcommand{\A}{\mathbb{A}}
\newcommand{\PP}{\mathbb{P}}
\newcommand{\E}{\mathbb{E}}
\newcommand{\eps}{\varepsilon}
\newcommand{\dd}{\,\mathrm{d}}
\begin{document}
\section*{UnsolvedMath ID: 10400122}
\textbf{Source problem code:} catalogue entry
\subsection{Definitions and mathematical statement}
Let \(M\) be a closed oriented three-manifold, and let \(\{c_0=0,c_1,\ldots,c_m\}\) be the source's finite set of values of the Chern--Simons functional on flat \(G\)-connections on \(M\). Let \(Z_k^G(M)\) be the quantum invariant at level \(k\), and put \(r=k+h^\vee\), where \(h^\vee\) is the indicated dual Coxeter number. The class of groups \(G\), normalization of the functional and invariant, and auxiliary framing/anomaly data are not specified in the excerpt.

The asymptotic expansion conjecture asserts that there exist, for each \(0\leq j\leq m\), numbers
\[d_j\in\Q,\quad\widetilde I_j\in\Q/\Z,\quad v_j\in\R_{>0},\quad a_j^e\in\C\quad(e\geq1),\]
independent of \(r\), such that, with \(d=\max_{0\leq j\leq m}d_j\),
\[Z_k^G(M)\sim\sum_{j=0}^m e^{2\pi i r c_j}r^{d_j}e^{\pi i\widetilde I_j/4}v_j
\left(1+\sum_{e=1}^{\infty}a_j^e r^{-e}\right)\qquad(r\to\infty).\]
Precisely, for every integer \(E\geq0\) there exists a constant \(C_E\geq0\), independent of \(r\), for which
\[\left|Z_k^G(M)-\sum_{j=0}^m e^{2\pi i r c_j}r^{d_j}e^{\pi i\widetilde I_j/4}v_j
\left(1+\sum_{e=1}^{E}a_j^e r^{-e}\right)\right|
\leq C_E r^{d-E-1}.\]
The source states this for all \(r=2,3,4,\ldots\); since \(r=k+h^\vee\), the meaning for levels not admissible for a chosen \(G\) requires clarification. Empty sums at \(E=0\) are zero. The constants and expansion data may depend on \(M,G\) and the fixed normalization, but \(C_E\) cannot depend on \(r\). This is an asymptotic expansion with truncation bounds, not a claim that the infinite series converges.

There is a literal phase ambiguity: \(e^{\pi i\widetilde I_j/4}\) is not defined by a class \(\widetilde I_j\in\Q/\Z\), since changing a representative by \(1\) changes the phase. The formula therefore retains the source's notation pending a specified lift or corrected quotient. If the \(c_j\) are classes modulo \(\Z\), their phases are well defined for integral \(r\), but that convention must likewise be stated. No lift or corrected phase group is silently selected.
\subsection{Short English statement}
Andersen's conjecture expresses quantum invariants at large level as finitely many Chern--Simons phases times rational powers and inverse-level asymptotic series, with a uniform bound after every truncation. The extracted phase quotient and level range need clarification.
\subsection{Sources}
\begin{itemize}
\item Local UnsolvedMath registry, permanent ID \(10400122\), exact source code \texttt{catalogue entry}.
\item J. E. Andersen, attribution and internal reference [6] in Ohtsuki, \emph{Problems on invariants of knots and 3-manifolds} (2002), Conjecture 7.7, PDF page 106: \url{https://msp.org/gtm/2002/04/gtm-2002-04-024s.pdf}. The bibliographic entry for [6] is absent. Registry access 2026-07-29; pdf-font-structure extraction.
\end{itemize}
\end{document}
